\documentclass{article}
% Source: Topic_04_Teacher_Edition.pdf, PDF pages 77-78 (printed pages 223A-223B): Topic Assessment Form A with the TE's marked answers, Form B with marked answers, Item Analysis for both forms. Cross-checked against the Realize downloads aga_24_a2_na_04_tta.docx (Form A) and aga_24_a2_na_04_ttb.docx (Form B), as text with the math dropped.
\usepackage{amsmath}
\usepackage{amssymb}
\newenvironment{assessment-meta}{}{}   % topic, title, forms, item count, printed pages, sources
\newenvironment{item-analysis}{}{}     % Savvas's Item / DOK / Skills Review and Practice table (same for both forms)
\newenvironment{ta-item}[2]{}{}        % #1 = item number, #2 = form (A or B)
\newcommand{\answer}[1]{}              % the TE's marked answer, inside a ta-item
\newcommand{\placeholder}[2]{}         % #1 = type (graph, table, figure, number-line), #2 = complete description
\newcommand{\uncertain}[2]{#1}         % #1 = best reading, #2 = alternatives and why
\newcommand{\te}[1]{}                  % transcriber's note (a discrepancy, a publisher misprint, a computed check)
\begin{document}
\begin{assessment-meta}
topic: 4
title: Rational Functions
forms: A and B (parallel item for item)
items: 19
printed: 223A-223B
sources: Topic_04_Teacher_Edition.pdf, PDF pages 77-78 (printed 223A-223B); Realize aga_24_a2_na_04_tta.docx and aga_24_a2_na_04_ttb.docx (text only, math dropped)
\end{assessment-meta}
\te{Student page headers read "4 Topic Assessment Form A" and "4 Topic Assessment Form B"; student footers read "enVision® Algebra 2 | 1of2 | Assessment Sourcebook" and "... 2of2 ..."; TE running head "TOPIC 4 Rational Functions", "TOPIC ASSESSMENT Form A" / "Form B"; TE footers "TOPIC 4 | 223A | TOPIC ASSESSMENT" and "TOPIC 4 | 223B | TOPIC ASSESSMENT". No running-head misprint on these pages. Every item, marked choice and boxed value was checked by computation; the only disagreement is Form B item 4 (see there).}
% Source page 78: 223B Item Analysis
\begin{item-analysis}
Item & DOK & Skills Review and Practice \\
1 & 2 & F74 \\
2 & 2 & F76 \\
3 & 2 & F76 \\
4 & 2 & A37 \\
5 & 1 & F75 \\
6 & 2 & A87 \\
7 & 3 & A36 \\
8 & 2 & A39 \\
9 & 2 & A87 \\
10 & 2 & A86 \\
11 & 2 & A87 \\
12 & 2 & F76 \\
13 & 3 & A87 \\
14 & 2 & F75 \\
15 & 2 & F74 \\
16 & 1 & F76 \\
17 & 2 & A87 \\
18 & 2 & F76 \\
19 & 2 & A74 \\
\end{item-analysis}
% Source page 77: 223A Topic Assessment Form A
\begin{ta-item}{1}{A}
The width, (y), of a rectangle with a fixed area varies inversely with its length, (x). The width is 3 inches when the length is 14 inches. Find the width when the length is 30 inches.
A. 0.71 inches
B. 4.48 inches
C. 6.43 inches
D. 1.4 inches
\answer{1A. D}
\end{ta-item}
\begin{ta-item}{2}{A}
What is the domain of the function (f(x)=\frac{x^2-x-2}{x^2-5x+6})?
A. All real numbers except 3
B. All real numbers except 2 and 3
C. All real numbers except -1 and 3
D. All real numbers except -1, 2, and 3
\answer{2A. B}
\end{ta-item}
\begin{ta-item}{3}{A}
What are the horizontal and vertical asymptotes of the graph of (y=\frac{x^2-3x-4}{1-3x-4x^2})?
A. (y=-\frac{1}{4}) and (x=-\frac{1}{4})
B. (y=-\frac{1}{4}) and (x=\frac{1}{4})
C. (y=\frac{1}{4}), (x=1), and (x=\frac{1}{4})
D. (y=-\frac{1}{4}), (x=-1), and (x=\frac{1}{4})
\answer{3A. B}
\end{ta-item}
\begin{ta-item}{4}{A}
Simplify (\frac{2x}{x-3}-\frac{-6}{3-x}). What are the restrictions on the domain of the expression?
[ ]
The domain is (\{x|x\neq) [ ](\}).
\answer{4A. 2; The domain is \{x|x \neq 3\}.}
\end{ta-item}
\begin{ta-item}{5}{A}
Describe the transformations needed to translate the graph of (y=\frac{1}{x}) to the graph of (y=\frac{1}{x+3}-4).
A. to the right 3 and up 4
B. to the left 3 and up 4
C. to the right 3 and down 4
D. to the left 3 and down 4
\answer{5A. D}
\end{ta-item}
\begin{ta-item}{6}{A}
Solve (\frac{10}{x^2-2x}=\frac{5}{x-2}-\frac{3}{x}).
A. (x=0)
B. (x=0) and 2
C. (x=2)
D. no solution
\answer{6A. D}
\end{ta-item}
\begin{ta-item}{7}{A}
Luke and Nora can peel 20 carrots in 6 min, working together. Luke can peel 5 carrots in 2 min, working alone. How many minutes would it take Nora to peel 10 carrots, working alone?
A. 15
B. 12
C. 8
D. 4
\answer{7A. B}
\end{ta-item}
\begin{ta-item}{8}{A}
What is the remainder when (5x^3-11x^2+5x+2) is divided by (5x-4)?
A. (-\frac{3}{5})
B. (\frac{38}{25})
C. (\frac{2}{5})
D. (-\frac{22}{5})
\answer{8A. B}
\end{ta-item}
\begin{ta-item}{9}{A}
Solve (\frac{3(1+x)}{4x}=1-\frac{1}{x}). (x=) [ ]
\answer{9A. 7}
\end{ta-item}
\begin{ta-item}{10}{A}
What is the sum (\frac{1}{x+3}+\frac{-6}{x^2-9})?
A. (\frac{-6}{x^2+x-12})
B. (\frac{1}{x^2+x+2})
C. (\frac{x-9}{x^2-9})
D. (\frac{-5}{x+3})
\answer{10A. C}
\end{ta-item}
\begin{ta-item}{11}{A}
Solve (\frac{1}{3}=\frac{x^2+2x-3}{3x^2-2x-1}).
A. (\frac{8}{3})
B. (-\frac{1}{3})
C. 1
D. no solution
\answer{11A. D}
\end{ta-item}
\begin{ta-item}{12}{A}
What are the horizontal and vertical asymptotes of the graph of (y=\frac{x+3}{x^4-81})?
A. (y=0; x=3)
B. (y=0; x=-3) and (x=3)
C. (y=1; x=9)
D. (y=1; x=-9)
\answer{12A. A}
\end{ta-item}
\begin{ta-item}{13}{A}
It takes 2 h for Faucet A to fill a tank, and it takes Faucet B 3 h. How many hours will it take the two faucets to fill the tank together?
A. 0.2
B. 1.2
C. 0.8
D. 1.8
\answer{13A. B}
\end{ta-item}
\begin{ta-item}{14}{A}
The graph of (xy=6) is translated up 3 units and to the right 2 units. Select all the possible equations for the translated graph.
A. (y=3+\frac{6}{x-2})
B. (y=\frac{3x}{x-2})
C. ((x-2)(y-3)=6)
D. (y=\frac{3}{x-2})
E. ((y-2)(x-3)=6)
\answer{14A. A, B, C}
\end{ta-item}
\begin{ta-item}{15}{A}
If (a=\frac{1}{x}+\frac{1}{y}), what is the value of (\frac{1}{a})?
A. (x+y)
B. (\frac{x+y}{xy})
C. (\frac{xy}{x+y})
D. (x^2+y^2)
\answer{15A. C}
\end{ta-item}
\begin{ta-item}{16}{A}
What are the horizontal and vertical asymptotes of the graph of (y=\frac{5x+1}{2x-5})?
A. (y=2.5; x=2.5)
B. (y=2.5; x=-2.5)
C. (y=0; x=2.5)
D. (y=0; x=-2.5)
\answer{16A. A}
\end{ta-item}
\begin{ta-item}{17}{A}
Solve (\frac{1}{x(x-3)}=\frac{2}{x-3}+\frac{3}{x}).
A. -2
B. 2
C. 1
D. -3
\answer{17A. B}
\end{ta-item}
\begin{ta-item}{18}{A}
Select all the functions whose graphs have a horizontal asymptote at (y=7).
A. (y=\frac{7}{x-2})
B. (y=\frac{7x^2+1}{x^2+7})
C. (y=7+\frac{1}{x})
D. (y=\frac{7x+3}{x^2+1})
E. (y=\frac{2x-1}{x^2-49})
\answer{18A. B, C}
\end{ta-item}
\begin{ta-item}{19}{A}
A rectangle has area (x^3+3x^2-18x) in.(^2) and width (x+6) in. What is the length in inches?
A. (x^2+3x)
B. (x^2+6x)
C. (x^2-3x)
D. (x^2-6x)
\answer{19A. C}
\end{ta-item}
% Source page 78: 223B Topic Assessment Form B
\begin{ta-item}{1}{B}
The width, (y), of a rectangle with a fixed area varies inversely with its length, (x). The width is 4 inches when the length is 18 inches. Find the width when the length is 40 inches.
A. 0.56 inches
B. 8.9 inches
C. 10.08 inches
D. 1.8 inches
\answer{1B. D}
\end{ta-item}
\begin{ta-item}{2}{B}
What is the domain of the function (f(x)=\frac{x^2-x-2}{x^4-81})?
A. All real numbers except 3
B. All real numbers except -1 and 3
C. All real numbers except -3 and 3
D. All real numbers except -3, 1, and 3
\answer{2B. C}
\end{ta-item}
\begin{ta-item}{3}{B}
What are the horizontal and vertical asymptotes of the graph of (y=\frac{x^2-3x-4}{3-x^2})?
A. (y=-1; x=\pm\sqrt{3})
B. (y=1; x=\pm\sqrt{3})
C. (y=-1; x=1) and (x=\sqrt{3})
D. (y=-1; x=1) and (x=-\sqrt{3})
\answer{3B. A}
\end{ta-item}
\begin{ta-item}{4}{B}
Simplify (\frac{1}{x-y}-\frac{-6}{y-x}). What are the any restrictions on the domain of the expression?
[ ]
The domain is (\{x|x\neq) [ ](\}).
\answer{4B. 2; The domain is \{x|x \neq 3\}.}
\te{computed: (\frac{1}{x-y}-\frac{-6}{y-x}=\frac{1}{x-y}-\frac{6}{x-y}=\frac{-5}{x-y}), restriction (x\neq y). The printed boxes (2 and 3) are Form A item 4's key carried over unchanged and do not fit this expression; publisher misprint in the key. "What are the any restrictions" is printed so on the student page and in the Realize Form B text.}
\end{ta-item}
\begin{ta-item}{5}{B}
Describe the transformations needed to translate the graph of (y=\frac{1}{x}) to the graph of (y=2+\frac{1}{x-5}).
A. to the left 5 and up 2
B. to the left 2 and down 5
C. to the right 2 and down 5
D. to the right 5 and up 2
\answer{5B. D}
\end{ta-item}
\begin{ta-item}{6}{B}
Solve (\frac{2x+4}{x^2+4x+3}=\frac{1}{x+1}+\frac{1}{x+3}).
A. no solution
B. (x=2)
C. (x=1) and 2
D. all values of (x), (x\neq-1) and (x\neq-3)
\answer{6B. D}
\end{ta-item}
\begin{ta-item}{7}{B}
Two robots can do a task in 5 min, working together. The first robot, working alone, can do the task in 15 min. How many minutes will it take the second robot, working alone, to do the task?
A. 10
B. 7.5
C. 5
D. 2
\answer{7B. B}
\end{ta-item}
\begin{ta-item}{8}{B}
What is the remainder when (4x^4-10x^2+2x+1) is divided by (4x-1)?
A. (-\frac{1}{2})
B. 4
C. (\frac{57}{64})
D. (-\frac{18}{37})
\answer{8B. C}
\end{ta-item}
\begin{ta-item}{9}{B}
Solve (\frac{2(1-x)}{3x}=1-\frac{3}{x}).
(x=) [ ]
\answer{9B. 2.2}
\end{ta-item}
\begin{ta-item}{10}{B}
What is the sum (\frac{1}{x-4}+\frac{-8}{x^2-16})?
A. (\frac{-8}{x^2+x-20})
B. (\frac{-7}{x+4})
C. (\frac{1}{x^2+x+2})
D. (\frac{1}{x+4})
\answer{10B. D}
\end{ta-item}
\begin{ta-item}{11}{B}
Solve (x=\frac{2x^2+x-7}{2x+8}).
A. 1
B. (\frac{1}{2})
C. -1
D. no solution
\answer{11B. C}
\end{ta-item}
\begin{ta-item}{12}{B}
What are the horizontal and vertical asymptotes of the graph of (y=\frac{x^4+3}{x^4+2x^2-3})?
A. (y=1; x=\pm1)
B. (y=1; x=\pm3)
C. (y=0; x=1)
D. (y=0; x=-1)
\answer{12B. A}
\te{The Realize Form B text gives choices A and B as "y = 1; x = +1" and "y = 1; x = +3"; the TE image prints (\pm) in both. The image rules (the text export dropped the (\pm) glyph).}
\end{ta-item}
\begin{ta-item}{13}{B}
It takes 4 h for Faucet A to fill a tank, and it takes Faucet B 6 h. How many hours will it take the two faucets to fill the tank together?
A. 1.4
B. 2.4
C. 2.0
D. 5.0
\answer{13B. B}
\end{ta-item}
\begin{ta-item}{14}{B}
The graph of (xy=6) is translated up 2 units and to the left 2 units. Select all the possible equations for the translated graph.
A. (y=2+\frac{6}{x+2})
B. (\frac{y}{2}=\frac{x+5}{x+2})
C. (y=\frac{2x+10}{x+2})
D. (y=\frac{6x+10}{x-2})
E. (y=4+\frac{4}{x-2})
\answer{14B. A, B, C}
\end{ta-item}
\begin{ta-item}{15}{B}
If (a=\frac{2}{x}+\frac{1}{y}), what is the value of (\frac{1}{a})?
A. (\frac{2y+x}{xy})
B. (\frac{x+2y}{2xy})
C. (\frac{xy}{x+2y})
D. (x+2y)
\answer{15B. C}
\end{ta-item}
\begin{ta-item}{16}{B}
What are the horizontal and vertical asymptotes of the graph of (y=\frac{-x+3}{x-8})?
A. (y=-1; x=8)
B. (y=-1; x=-8)
C. (y=1; x=8)
D. (y=1; x=-8)
\answer{16B. A}
\end{ta-item}
\begin{ta-item}{17}{B}
Solve (\frac{x^2+x-4}{x-2}=x-\frac{1}{x-2}).
A. -2
B. 1
C. -1
D. 3
\answer{17B. B}
\end{ta-item}
\begin{ta-item}{18}{B}
Select all the functions whose graphs have a horizontal asymptote at (y=\frac{2}{3}).
A. (y=\frac{2}{3x-1})
B. (y=\frac{2x^2+1}{3x^2-2})
C. (y=\frac{2}{3}+\frac{1}{x})
D. (y=\frac{2x-3}{3x^2+1})
E. (y=3+\frac{3}{2x})
\answer{18B. B, C}
\end{ta-item}
\begin{ta-item}{19}{B}
A rectangle has area (x^3-15x-4) cm(^2) and width (x-4) cm. What is the length in centimeters?
A. (x^2+4x)
B. (x^2-4x+2)
C. (x^2+4x+1)
D. (x^2+1)
\answer{19B. C}
\end{ta-item}
\end{document}
